\section{Learning Objectives}

After completing this chapter, readers will be able to:
\begin{itemize}
    \item Apply classical hypothesis testing frameworks to A/B tests
    \item Select appropriate statistical tests for different metric types
    \item Check assumptions and diagnose violations
    \item Implement non-parametric alternatives when assumptions fail
    \item Interpret test results correctly and avoid common pitfalls
    \item Use decision frameworks for systematic test selection
\end{itemize}

\section{Classical Hypothesis Testing Framework}

Hypothesis testing provides a formal procedure for making decisions under uncertainty. While introduced earlier, this chapter provides rigorous treatment and practical implementation.

\subsection{Structure of a Hypothesis Test}

\begin{enumerate}
    \item \textbf{State hypotheses:}
    \begin{itemize}
        \item Null hypothesis $H_0$: No effect or no difference
        \item Alternative hypothesis $H_1$: An effect exists
    \end{itemize}

    \item \textbf{Choose significance level:} $\alpha$ (typically 0.05 or 0.01)

    \item \textbf{Calculate test statistic:} Function of data measuring evidence against $H_0$

    \item \textbf{Determine p-value:} Probability of observing data this extreme under $H_0$

    \item \textbf{Make decision:}
    \begin{itemize}
        \item Reject $H_0$ if p-value $< \alpha$
        \item Fail to reject $H_0$ otherwise
    \end{itemize}

    \item \textbf{Interpret in context:} Translate statistical conclusion to practical meaning
\end{enumerate}

\subsection{One-Sided vs. Two-Sided Tests}

\textbf{Two-sided test:}
\begin{equation}
H_0: \theta_1 = \theta_2 \quad \text{vs.} \quad H_1: \theta_1 \neq \theta_2
\end{equation}

Tests for any difference (positive or negative). Most common in A/B testing.

\textbf{One-sided test:}
\begin{equation}
H_0: \theta_1 \geq \theta_2 \quad \text{vs.} \quad H_1: \theta_1 < \theta_2
\end{equation}

Tests specifically for improvement (or degradation). More powerful but only detects effects in specified direction.

\begin{theorem}[Relationship Between One-Sided and Two-Sided Tests]
For the same data and test statistic, the one-sided p-value equals half the two-sided p-value (when the effect is in the hypothesized direction).

If $p_{\text{two-sided}} = 2P(|Z| > |z_{\text{obs}}|)$, then:
\begin{equation}
p_{\text{one-sided}} = P(Z > z_{\text{obs}}) = \frac{p_{\text{two-sided}}}{2}
\end{equation}
\end{theorem}

\textbf{When to use one-sided tests:}
\begin{itemize}
    \item Only care about improvements (would not implement if treatment is worse)
    \item Regulatory or business constraints dictate directionality
    \item Pre-registered directional hypothesis
    \item Strong theoretical reason to expect one direction only
\end{itemize}

\textbf{Caution:} One-sided tests are often misused. Default to two-sided unless strong justification exists.

\section{Tests for Proportions}

Proportion tests are fundamental for A/B testing of conversion rates, click-through rates, and other binary outcomes.

\subsection{Two-Sample Z-Test for Proportions}

For control proportion $\hat{p}_1 = x_1/n_1$ and treatment proportion $\hat{p}_2 = x_2/n_2$:

\begin{equation}
H_0: p_1 = p_2 \quad \text{vs.} \quad H_1: p_1 \neq p_2
\end{equation}

Under $H_0$, we use the pooled proportion:

\begin{equation}
\hat{p}_{\text{pooled}} = \frac{x_1 + x_2}{n_1 + n_2}
\end{equation}

Test statistic:

\begin{equation}
Z = \frac{\hat{p}_2 - \hat{p}_1}{\sqrt{\hat{p}_{\text{pooled}}(1 - \hat{p}_{\text{pooled}})(1/n_1 + 1/n_2)}}
\end{equation}

Under $H_0$: $Z \sim \mathcal{N}(0, 1)$ approximately (by CLT for large samples).

P-value (two-sided):
\begin{equation}
p = 2(1 - \Phi(|Z|))
\end{equation}
where $\Phi$ is the standard normal CDF.

Decision rule: Reject $H_0$ if $|Z| > z_{\alpha/2}$ (equivalently, if $p < \alpha$).

\subsection{Confidence Interval for Difference in Proportions}

For confidence interval, use unpooled variance:

\begin{equation}
\text{SE}_{\text{unpooled}} = \sqrt{\frac{\hat{p}_1(1-\hat{p}_1)}{n_1} + \frac{\hat{p}_2(1-\hat{p}_2)}{n_2}}
\end{equation}

$(1-\alpha)$ confidence interval:
\begin{equation}
(\hat{p}_2 - \hat{p}_1) \pm z_{\alpha/2} \cdot \text{SE}_{\text{unpooled}}
\end{equation}

\begin{example}[Conversion Rate Test]
\begin{align*}
\text{Control: } & n_1 = 2000, x_1 = 160, \hat{p}_1 = 0.080 \\
\text{Treatment: } & n_2 = 2000, x_2 = 190, \hat{p}_2 = 0.095
\end{align*}

Pooled proportion:
\begin{equation}
\hat{p}_{\text{pooled}} = \frac{160 + 190}{2000 + 2000} = \frac{350}{4000} = 0.0875
\end{equation}

Standard error:
\begin{equation}
\text{SE} = \sqrt{0.0875 \times 0.9125 \times (1/2000 + 1/2000)} = 0.00894
\end{equation}

Test statistic:
\begin{equation}
Z = \frac{0.095 - 0.080}{0.00894} = 1.68
\end{equation}

P-value: $p = 2(1 - \Phi(1.68)) = 2(1 - 0.9535) = 0.093$

At $\alpha = 0.05$, we fail to reject $H_0$. The observed difference is not statistically significant.
\end{example}

\lstset{language=Python, caption=Comprehensive Z-Test Implementation for Proportions}
\begin{lstlisting}
import numpy as np
from scipy import stats
import matplotlib.pyplot as plt
from typing import Dict, Tuple
import pandas as pd

def z_test_proportions(
    x1: int,
    n1: int,
    x2: int,
    n2: int,
    alpha: float = 0.05,
    alternative: str = 'two-sided'
) -> Dict:
    """
    Perform z-test for comparing two proportions.

    Parameters:
    -----------
    x1, x2 : int
        Number of successes in each group
    n1, n2 : int
        Sample sizes
    alpha : float
        Significance level
    alternative : str
        'two-sided', 'greater', or 'less'

    Returns:
    --------
    dict : Test results including statistic, p-value, CI, decision
    """
    # Sample proportions
    p1_hat = x1 / n1
    p2_hat = x2 / n2
    diff = p2_hat - p1_hat

    # Pooled proportion under null hypothesis
    p_pooled = (x1 + x2) / (n1 + n2)

    # Standard error under null (pooled)
    se_pooled = np.sqrt(p_pooled * (1 - p_pooled) * (1/n1 + 1/n2))

    # Test statistic
    z_stat = diff / se_pooled

    # P-value
    if alternative == 'two-sided':
        p_value = 2 * (1 - stats.norm.cdf(abs(z_stat)))
    elif alternative == 'greater':
        p_value = 1 - stats.norm.cdf(z_stat)
    elif alternative == 'less':
        p_value = stats.norm.cdf(z_stat)
    else:
        raise ValueError("alternative must be 'two-sided', 'greater', or 'less'")

    # Decision
    reject_null = p_value < alpha

    # Confidence interval (using unpooled variance)
    se_unpooled = np.sqrt(p1_hat*(1-p1_hat)/n1 + p2_hat*(1-p2_hat)/n2)

    if alternative == 'two-sided':
        z_crit = stats.norm.ppf(1 - alpha/2)
        ci_lower = diff - z_crit * se_unpooled
        ci_upper = diff + z_crit * se_unpooled
    elif alternative == 'greater':
        z_crit = stats.norm.ppf(1 - alpha)
        ci_lower = diff - z_crit * se_unpooled
        ci_upper = np.inf
    else:  # less
        z_crit = stats.norm.ppf(1 - alpha)
        ci_lower = -np.inf
        ci_upper = diff + z_crit * se_unpooled

    # Effect size (relative lift)
    relative_lift = diff / p1_hat if p1_hat > 0 else np.nan

    return {
        'test': 'Z-test for proportions',
        'p1_hat': p1_hat,
        'p2_hat': p2_hat,
        'difference': diff,
        'relative_lift': relative_lift,
        'z_statistic': z_stat,
        'p_value': p_value,
        'alpha': alpha,
        'reject_null': reject_null,
        'ci_lower': ci_lower,
        'ci_upper': ci_upper,
        'se_pooled': se_pooled,
        'se_unpooled': se_unpooled,
        'alternative': alternative,
        'decision': 'Reject H0' if reject_null else 'Fail to reject H0'
    }

def visualize_proportion_test(result: Dict, x1: int, n1: int, x2: int, n2: int):
    """Visualize proportion test results."""
    fig, axes = plt.subplots(1, 3, figsize=(15, 5))

    # 1. Proportions with CI
    ax = axes[0]
    props = [result['p1_hat'], result['p2_hat']]
    labels = ['Control', 'Treatment']
    colors = ['C0', 'C1']

    # Standard errors for individual proportions
    se1 = np.sqrt(result['p1_hat'] * (1 - result['p1_hat']) / n1)
    se2 = np.sqrt(result['p2_hat'] * (1 - result['p2_hat']) / n2)
    errors = [1.96 * se1, 1.96 * se2]

    x_pos = np.arange(2)
    bars = ax.bar(x_pos, props, yerr=errors, capsize=5, alpha=0.7, color=colors)
    ax.set_xticks(x_pos)
    ax.set_xticklabels(labels)
    ax.set_ylabel('Conversion Rate')
    ax.set_title('Conversion Rates with 95% CI')
    ax.grid(axis='y', alpha=0.3)

    # Add values on bars
    for i, (bar, prop) in enumerate(zip(bars, props)):
        height = bar.get_height()
        ax.text(bar.get_x() + bar.get_width()/2., height + errors[i] + 0.01,
                f'{prop:.3f}', ha='center', va='bottom', fontweight='bold')

    # 2. Difference with CI
    ax = axes[1]
    diff = result['difference']
    ci_lower = result['ci_lower']
    ci_upper = result['ci_upper']

    ax.barh([0], [diff], xerr=[[diff - ci_lower], [ci_upper - diff]],
            capsize=5, alpha=0.7, color='green' if diff > 0 else 'red')
    ax.axvline(x=0, color='black', linestyle='--', linewidth=1)
    ax.set_yticks([])
    ax.set_xlabel('Difference in Proportions')
    ax.set_title(f'Difference: {diff:.4f} [{ci_lower:.4f}, {ci_upper:.4f}]')
    ax.grid(axis='x', alpha=0.3)

    # Add significance marker
    if result['reject_null']:
        ax.text(diff, 0.3, 'Significant', ha='center', fontsize=12,
                fontweight='bold', color='darkgreen')
    else:
        ax.text(diff, 0.3, 'Not Significant', ha='center', fontsize=12,
                fontweight='bold', color='gray')

    # 3. Null distribution with test statistic
    ax = axes[2]
    x = np.linspace(-4, 4, 1000)
    y = stats.norm.pdf(x, 0, 1)
    ax.plot(x, y, 'b-', linewidth=2, label='Null distribution N(0,1)')
    ax.axvline(x=result['z_statistic'], color='red', linestyle='--',
               linewidth=2, label=f'Observed z = {result["z_statistic"]:.2f}')

    # Shade rejection region
    if result['alternative'] == 'two-sided':
        z_crit = stats.norm.ppf(1 - result['alpha']/2)
        x_left = x[x <= -z_crit]
        x_right = x[x >= z_crit]
        ax.fill_between(x_left, 0, stats.norm.pdf(x_left, 0, 1),
                        alpha=0.3, color='red')
        ax.fill_between(x_right, 0, stats.norm.pdf(x_right, 0, 1),
                        alpha=0.3, color='red', label=f'Rejection region (α={result["alpha"]})')
        ax.axvline(x=-z_crit, color='red', linestyle=':', alpha=0.5)
        ax.axvline(x=z_crit, color='red', linestyle=':', alpha=0.5)

    ax.set_xlabel('Z-statistic')
    ax.set_ylabel('Density')
    ax.set_title(f'Null Distribution (p-value = {result["p_value"]:.4f})')
    ax.legend()
    ax.grid(alpha=0.3)

    plt.tight_layout()
    return fig

# Example usage
print("=== Z-Test for Proportions ===\n")

# Example 1: Non-significant result
result1 = z_test_proportions(x1=160, n1=2000, x2=190, n2=2000)

print("Example 1: Marginal effect")
print(f"Control:    {result1['p1_hat']:.3f} ({160}/{2000})")
print(f"Treatment:  {result1['p2_hat']:.3f} ({190}/{2000})")
print(f"Difference: {result1['difference']:.4f}")
print(f"Relative lift: {result1['relative_lift']*100:.1f}%")
print(f"Z-statistic: {result1['z_statistic']:.3f}")
print(f"P-value: {result1['p_value']:.4f}")
print(f"95% CI: [{result1['ci_lower']:.4f}, {result1['ci_upper']:.4f}]")
print(f"Decision: {result1['decision']}\n")

# Example 2: Significant result
result2 = z_test_proportions(x1=200, n1=2000, x2=260, n2=2000)

print("Example 2: Strong effect")
print(f"Control:    {result2['p1_hat']:.3f} ({200}/{2000})")
print(f"Treatment:  {result2['p2_hat']:.3f} ({260}/{2000})")
print(f"Difference: {result2['difference']:.4f}")
print(f"Relative lift: {result2['relative_lift']*100:.1f}%")
print(f"Z-statistic: {result2['z_statistic']:.3f}")
print(f"P-value: {result2['p_value']:.4e}")
print(f"95% CI: [{result2['ci_lower']:.4f}, {result2['ci_upper']:.4f}]")
print(f"Decision: {result2['decision']}\n")

# Example 3: One-sided test
result3 = z_test_proportions(x1=200, n1=2000, x2=260, n2=2000,
                             alternative='greater')

print("Example 3: One-sided test (greater)")
print(f"Z-statistic: {result3['z_statistic']:.3f}")
print(f"P-value (one-sided): {result3['p_value']:.4e}")
print(f"P-value (two-sided): {result2['p_value']:.4e}")
print(f"Ratio: {result2['p_value'] / result3['p_value']:.1f}\n")

# Visualize
fig1 = visualize_proportion_test(result1, 160, 2000, 190, 2000)
plt.savefig('proportion_test_marginal.pdf')

fig2 = visualize_proportion_test(result2, 200, 2000, 260, 2000)
plt.savefig('proportion_test_strong.pdf')

plt.show()
\end{lstlisting}

\subsection{Assumptions for Z-Test}

The z-test for proportions relies on:

\begin{enumerate}
    \item \textbf{Independence:} Observations are independent within and between groups

    \item \textbf{Random sampling:} Groups formed by randomization

    \item \textbf{Large samples:} Normal approximation valid

    Rule of thumb: $n_i p_i > 5$ and $n_i(1-p_i) > 5$ for $i = 1, 2$.

    More conservative: $n_i p_i > 10$ and $n_i(1-p_i) > 10$.
\end{enumerate}

When assumptions are violated, consider:
\begin{itemize}
    \item Fisher's exact test for small samples
    \item Permutation tests (robust to distributional assumptions)
    \item Adjustments for clustering or correlation
\end{itemize}

\section{Tests for Means}

When comparing continuous outcomes (revenue, time on site, engagement metrics), we use t-tests.

\subsection{Two-Sample T-Test}

For means $\bar{X}_1$ and $\bar{X}_2$ from groups with variances $s_1^2$ and $s_2^2$:

\begin{equation}
H_0: \mu_1 = \mu_2 \quad \text{vs.} \quad H_1: \mu_1 \neq \mu_2
\end{equation}

\textbf{Pooled (Student's) t-test} (assumes equal variances):

Pooled variance:
\begin{equation}
s_p^2 = \frac{(n_1 - 1)s_1^2 + (n_2 - 1)s_2^2}{n_1 + n_2 - 2}
\end{equation}

Test statistic:
\begin{equation}
t = \frac{\bar{X}_2 - \bar{X}_1}{s_p \sqrt{1/n_1 + 1/n_2}}
\end{equation}

Distribution under $H_0$: $t \sim t_{n_1 + n_2 - 2}$

\textbf{Welch's t-test} (unequal variances):

Test statistic:
\begin{equation}
t = \frac{\bar{X}_2 - \bar{X}_1}{\sqrt{s_1^2/n_1 + s_2^2/n_2}}
\end{equation}

Degrees of freedom (Welch-Satterthwaite approximation):
\begin{equation}
\nu = \frac{(s_1^2/n_1 + s_2^2/n_2)^2}{\frac{(s_1^2/n_1)^2}{n_1-1} + \frac{(s_2^2/n_2)^2}{n_2-1}}
\end{equation}

\begin{theorem}[Robustness of Welch's t-test]
Welch's t-test maintains correct Type I error rate even when:
\begin{itemize}
    \item Variances are unequal
    \item Sample sizes are unequal
    \item Both conditions hold simultaneously
\end{itemize}

The pooled t-test can have inflated Type I error when $\sigma_1 \neq \sigma_2$ and $n_1 \neq n_2$.
\end{theorem}

\textbf{Practical recommendation:} Default to Welch's t-test. It performs nearly as well as pooled t-test when variances are equal, and much better when they're not.

\lstset{language=Python, caption=Comprehensive T-Test Implementation}
\begin{lstlisting}
import numpy as np
from scipy import stats
import matplotlib.pyplot as plt
from typing import Dict
import pandas as pd

def t_test_means(
    data1: np.ndarray,
    data2: np.ndarray,
    alpha: float = 0.05,
    equal_var: bool = False,
    alternative: str = 'two-sided'
) -> Dict:
    """
    Perform t-test for comparing two means.

    Parameters:
    -----------
    data1, data2 : np.ndarray
        Data for each group
    alpha : float
        Significance level
    equal_var : bool
        If True, use pooled t-test. If False, use Welch's t-test
    alternative : str
        'two-sided', 'greater', or 'less'

    Returns:
    --------
    dict : Test results
    """
    n1, n2 = len(data1), len(data2)
    mean1, mean2 = np.mean(data1), np.mean(data2)
    var1, var2 = np.var(data1, ddof=1), np.var(data2, ddof=1)
    std1, std2 = np.std(data1, ddof=1), np.std(data2, ddof=1)

    diff = mean2 - mean1

    if equal_var:
        # Pooled t-test
        pooled_var = ((n1 - 1) * var1 + (n2 - 1) * var2) / (n1 + n2 - 2)
        se = np.sqrt(pooled_var * (1/n1 + 1/n2))
        df = n1 + n2 - 2
        test_type = "Pooled t-test"
    else:
        # Welch's t-test
        se = np.sqrt(var1/n1 + var2/n2)
        # Welch-Satterthwaite degrees of freedom
        df = (var1/n1 + var2/n2)**2 / (
            (var1/n1)**2/(n1-1) + (var2/n2)**2/(n2-1)
        )
        test_type = "Welch's t-test"

    # Test statistic
    t_stat = diff / se

    # P-value
    if alternative == 'two-sided':
        p_value = 2 * (1 - stats.t.cdf(abs(t_stat), df))
    elif alternative == 'greater':
        p_value = 1 - stats.t.cdf(t_stat, df)
    elif alternative == 'less':
        p_value = stats.t.cdf(t_stat, df)
    else:
        raise ValueError("alternative must be 'two-sided', 'greater', or 'less'")

    # Decision
    reject_null = p_value < alpha

    # Confidence interval
    if alternative == 'two-sided':
        t_crit = stats.t.ppf(1 - alpha/2, df)
        ci_lower = diff - t_crit * se
        ci_upper = diff + t_crit * se
    elif alternative == 'greater':
        t_crit = stats.t.ppf(1 - alpha, df)
        ci_lower = diff - t_crit * se
        ci_upper = np.inf
    else:  # less
        t_crit = stats.t.ppf(1 - alpha, df)
        ci_lower = -np.inf
        ci_upper = diff + t_crit * se

    # Effect size (Cohen's d)
    if equal_var:
        pooled_std = np.sqrt(pooled_var)
    else:
        # Use pooled std even for Welch's when calculating Cohen's d
        pooled_std = np.sqrt(((n1-1)*var1 + (n2-1)*var2) / (n1 + n2 - 2))

    cohens_d = diff / pooled_std

    # Relative effect (if mean1 > 0)
    relative_diff = diff / mean1 if mean1 != 0 else np.nan

    return {
        'test': test_type,
        'n1': n1,
        'n2': n2,
        'mean1': mean1,
        'mean2': mean2,
        'std1': std1,
        'std2': std2,
        'difference': diff,
        'relative_diff': relative_diff,
        't_statistic': t_stat,
        'df': df,
        'p_value': p_value,
        'alpha': alpha,
        'reject_null': reject_null,
        'ci_lower': ci_lower,
        'ci_upper': ci_upper,
        'se': se,
        'cohens_d': cohens_d,
        'alternative': alternative,
        'decision': 'Reject H0' if reject_null else 'Fail to reject H0'
    }

def compare_t_test_variants(data1: np.ndarray, data2: np.ndarray, alpha: float = 0.05):
    """Compare pooled vs. Welch's t-test."""

    pooled = t_test_means(data1, data2, alpha=alpha, equal_var=True)
    welch = t_test_means(data1, data2, alpha=alpha, equal_var=False)

    print("=== Comparison: Pooled vs. Welch's t-test ===\n")

    print(f"Sample sizes: n1={len(data1)}, n2={len(data2)}")
    print(f"Means: {np.mean(data1):.2f}, {np.mean(data2):.2f}")
    print(f"Std devs: {np.std(data1, ddof=1):.2f}, {np.std(data2, ddof=1):.2f}")
    print(f"Variance ratio: {np.var(data2, ddof=1) / np.var(data1, ddof=1):.2f}\n")

    comparison = pd.DataFrame({
        'Pooled t-test': [
            pooled['t_statistic'],
            pooled['df'],
            pooled['p_value'],
            pooled['ci_lower'],
            pooled['ci_upper'],
            pooled['decision']
        ],
        'Welch\'s t-test': [
            welch['t_statistic'],
            welch['df'],
            welch['p_value'],
            welch['ci_lower'],
            welch['ci_upper'],
            welch['decision']
        ]
    }, index=['t-statistic', 'df', 'p-value', 'CI lower', 'CI upper', 'Decision'])

    print(comparison)
    print()

    return pooled, welch

def visualize_t_test(result: Dict, data1: np.ndarray, data2: np.ndarray):
    """Visualize t-test results."""
    fig, axes = plt.subplots(2, 2, figsize=(14, 10))

    # 1. Distributions
    ax = axes[0, 0]
    ax.hist(data1, bins=30, alpha=0.6, label='Control', density=True, color='C0')
    ax.hist(data2, bins=30, alpha=0.6, label='Treatment', density=True, color='C1')
    ax.axvline(result['mean1'], color='C0', linestyle='--', linewidth=2,
              label=f'Control mean = {result["mean1"]:.2f}')
    ax.axvline(result['mean2'], color='C1', linestyle='--', linewidth=2,
              label=f'Treatment mean = {result["mean2"]:.2f}')
    ax.set_xlabel('Value')
    ax.set_ylabel('Density')
    ax.set_title('Data Distributions')
    ax.legend()
    ax.grid(alpha=0.3)

    # 2. Means with CI
    ax = axes[0, 1]
    means = [result['mean1'], result['mean2']]
    sems = [result['std1']/np.sqrt(result['n1']),
            result['std2']/np.sqrt(result['n2'])]
    labels = ['Control', 'Treatment']

    x_pos = np.arange(2)
    bars = ax.bar(x_pos, means, yerr=[1.96*s for s in sems],
                  capsize=5, alpha=0.7, color=['C0', 'C1'])
    ax.set_xticks(x_pos)
    ax.set_xticklabels(labels)
    ax.set_ylabel('Mean Value')
    ax.set_title('Means with 95% CI')
    ax.grid(axis='y', alpha=0.3)

    # Add values
    for i, (bar, mean) in enumerate(zip(bars, means)):
        height = bar.get_height()
        ax.text(bar.get_x() + bar.get_width()/2., height + sems[i]*2,
                f'{mean:.2f}', ha='center', va='bottom', fontweight='bold')

    # 3. Difference with CI
    ax = axes[1, 0]
    diff = result['difference']
    ci_lower = result['ci_lower']
    ci_upper = result['ci_upper']

    ax.barh([0], [diff], xerr=[[diff - ci_lower], [ci_upper - diff]],
            capsize=5, alpha=0.7, color='green' if diff > 0 else 'red')
    ax.axvline(x=0, color='black', linestyle='--', linewidth=1)
    ax.set_yticks([])
    ax.set_xlabel('Difference in Means')
    ax.set_title(f'{result["test"]}\nDifference: {diff:.2f} '
                f'[{ci_lower:.2f}, {ci_upper:.2f}]')
    ax.grid(axis='x', alpha=0.3)

    # Add significance and effect size
    if result['reject_null']:
        sig_text = 'Significant'
        color = 'darkgreen'
    else:
        sig_text = 'Not Significant'
        color = 'gray'

    ax.text(diff, 0.3, f'{sig_text}\nCohen\'s d = {result["cohens_d"]:.3f}',
           ha='center', fontsize=10, fontweight='bold', color=color)

    # 4. T-distribution with test statistic
    ax = axes[1, 1]
    df = result['df']
    x = np.linspace(-5, 5, 1000)
    y = stats.t.pdf(x, df)
    ax.plot(x, y, 'b-', linewidth=2, label=f't-distribution (df={df:.1f})')
    ax.axvline(x=result['t_statistic'], color='red', linestyle='--',
               linewidth=2, label=f'Observed t = {result["t_statistic"]:.2f}')

    # Shade rejection region for two-sided
    if result['alternative'] == 'two-sided':
        t_crit = stats.t.ppf(1 - result['alpha']/2, df)
        x_left = x[x <= -t_crit]
        x_right = x[x >= t_crit]
        ax.fill_between(x_left, 0, stats.t.pdf(x_left, df),
                        alpha=0.3, color='red')
        ax.fill_between(x_right, 0, stats.t.pdf(x_right, df),
                        alpha=0.3, color='red',
                        label=f'Rejection region (α={result["alpha"]})')
        ax.axvline(x=-t_crit, color='red', linestyle=':', alpha=0.5)
        ax.axvline(x=t_crit, color='red', linestyle=':', alpha=0.5)

    ax.set_xlabel('t-statistic')
    ax.set_ylabel('Density')
    ax.set_title(f't-Distribution (p-value = {result["p_value"]:.4f})')
    ax.legend()
    ax.grid(alpha=0.3)

    plt.tight_layout()
    return fig

# Example usage
print("=== T-Test for Means ===\n")

# Generate data
np.random.seed(42)

# Example 1: Equal variances
data1_eq = np.random.normal(45, 10, 500)
data2_eq = np.random.normal(48, 10, 500)

pooled_eq, welch_eq = compare_t_test_variants(data1_eq, data2_eq)

# Example 2: Unequal variances
data1_uneq = np.random.normal(45, 10, 500)
data2_uneq = np.random.normal(48, 20, 500)

print("\n=== Example 2: Unequal Variances ===\n")
pooled_uneq, welch_uneq = compare_t_test_variants(data1_uneq, data2_uneq)

# Visualize
fig1 = visualize_t_test(welch_eq, data1_eq, data2_eq)
plt.savefig('t_test_equal_var.pdf')

fig2 = visualize_t_test(welch_uneq, data1_uneq, data2_uneq)
plt.savefig('t_test_unequal_var.pdf')

plt.show()
\end{lstlisting}

\section{Chi-Squared Test for Categorical Data}

The chi-squared test assesses association between categorical variables, useful for multi-category outcomes.

\subsection{Pearson's Chi-Squared Test}

For $r \times c$ contingency table with observed counts $O_{ij}$:

\begin{equation}
H_0: \text{Row and column variables are independent}
\end{equation}

Expected counts under independence:
\begin{equation}
E_{ij} = \frac{(\text{row } i \text{ total}) \times (\text{column } j \text{ total})}{n}
\end{equation}

Test statistic:
\begin{equation}
\chi^2 = \sum_{i=1}^r \sum_{j=1}^c \frac{(O_{ij} - E_{ij})^2}{E_{ij}}
\end{equation}

Under $H_0$: $\chi^2 \sim \chi^2_{(r-1)(c-1)}$ approximately.

\begin{proposition}[Equivalence for 2×2 Tables]
For $2 \times 2$ contingency tables, $\chi^2 = Z^2$ where $Z$ is the z-test statistic for proportions. The tests yield identical p-values for two-sided alternatives.
\end{proposition}

\textbf{Assumptions:}
\begin{itemize}
    \item Independent observations
    \item Expected cell counts $E_{ij} \geq 5$ for all cells
    \item If some $E_{ij} < 5$: use Fisher's exact test or combine categories
\end{itemize}

\subsection{Fisher's Exact Test}

For small samples or when chi-squared assumptions fail, Fisher's exact test calculates exact p-values.

Particularly useful for $2 \times 2$ tables with small expected counts.

\lstset{language=Python, caption=Chi-Squared and Fisher's Exact Test Implementation}
\begin{lstlisting}
from scipy.stats import chi2_contingency, fisher_exact
import numpy as np
import pandas as pd
import matplotlib.pyplot as plt
from typing import Dict, Tuple

def chi_square_test(
    table: np.ndarray,
    alpha: float = 0.05
) -> Dict:
    """
    Perform chi-squared test of independence.

    Parameters:
    -----------
    table : np.ndarray
        Contingency table (rows × columns)
    alpha : float
        Significance level

    Returns:
    --------
    dict : Test results
    """
    chi2_stat, p_value, dof, expected = chi2_contingency(table)

    reject_null = p_value < alpha

    # Check assumption: expected counts >= 5
    min_expected = expected.min()
    assumption_met = min_expected >= 5

    # Cramer's V effect size
    n = table.sum()
    min_dim = min(table.shape[0] - 1, table.shape[1] - 1)
    cramers_v = np.sqrt(chi2_stat / (n * min_dim))

    return {
        'test': 'Chi-squared test',
        'chi2_statistic': chi2_stat,
        'p_value': p_value,
        'df': dof,
        'expected_counts': expected,
        'min_expected': min_expected,
        'assumption_met': assumption_met,
        'cramers_v': cramers_v,
        'alpha': alpha,
        'reject_null': reject_null,
        'decision': 'Reject H0' if reject_null else 'Fail to reject H0'
    }

def fishers_test_2x2(
    table: np.ndarray,
    alpha: float = 0.05,
    alternative: str = 'two-sided'
) -> Dict:
    """
    Perform Fisher's exact test for 2×2 table.

    Parameters:
    -----------
    table : np.ndarray
        2×2 contingency table
    alpha : float
        Significance level
    alternative : str
        'two-sided', 'greater', or 'less'

    Returns:
    --------
    dict : Test results
    """
    if table.shape != (2, 2):
        raise ValueError("Fisher's exact test requires a 2×2 table")

    oddsratio, p_value = fisher_exact(table, alternative=alternative)

    reject_null = p_value < alpha

    return {
        'test': 'Fisher\'s exact test',
        'odds_ratio': oddsratio,
        'p_value': p_value,
        'alpha': alpha,
        'reject_null': reject_null,
        'alternative': alternative,
        'decision': 'Reject H0' if reject_null else 'Fail to reject H0'
    }

def compare_chi_square_fisher(table: np.ndarray, alpha: float = 0.05):
    """Compare chi-squared test with Fisher's exact test."""

    print("=== Comparison: Chi-Squared vs. Fisher's Exact Test ===\n")

    print("Contingency table:")
    print(table)
    print()

    # Chi-squared test
    chi2_result = chi_square_test(table, alpha)

    print("Chi-Squared Test:")
    print(f"  Chi-squared statistic: {chi2_result['chi2_statistic']:.3f}")
    print(f"  Degrees of freedom: {chi2_result['df']}")
    print(f"  P-value: {chi2_result['p_value']:.4f}")
    print(f"  Cram\'er's V: {chi2_result['cramers_v']:.3f}")
    print(f"  Min expected count: {chi2_result['min_expected']:.2f}")
    print(f"  Assumption met (E>=5): {chi2_result['assumption_met']}")
    print(f"  Decision: {chi2_result['decision']}\n")

    print("Expected frequencies:")
    print(chi2_result['expected_counts'].round(2))
    print()

    # Fisher's exact test (if 2×2)
    if table.shape == (2, 2):
        fisher_result = fishers_test_2x2(table, alpha)

        print("Fisher's Exact Test:")
        print(f"  Odds ratio: {fisher_result['odds_ratio']:.3f}")
        print(f"  P-value: {fisher_result['p_value']:.4f}")
        print(f"  Decision: {fisher_result['decision']}\n")

        # Compare p-values
        print(f"P-value difference: {abs(chi2_result['p_value'] - fisher_result['p_value']):.4f}")

        if not chi2_result['assumption_met']:
            print("\nNote: Expected counts < 5. Fisher's exact test is more reliable.")

        return chi2_result, fisher_result
    else:
        return chi2_result, None

def visualize_contingency_table(table: np.ndarray, result: Dict,
                                row_labels=None, col_labels=None):
    """Visualize contingency table analysis."""
    fig, axes = plt.subplots(1, 3, figsize=(15, 5))

    if row_labels is None:
        row_labels = [f'Row {i}' for i in range(table.shape[0])]
    if col_labels is None:
        col_labels = [f'Col {j}' for j in range(table.shape[1])]

    # 1. Observed counts heatmap
    ax = axes[0]
    im1 = ax.imshow(table, cmap='YlOrRd', aspect='auto')
    ax.set_xticks(range(table.shape[1]))
    ax.set_yticks(range(table.shape[0]))
    ax.set_xticklabels(col_labels)
    ax.set_yticklabels(row_labels)
    ax.set_title('Observed Counts')

    # Add text annotations
    for i in range(table.shape[0]):
        for j in range(table.shape[1]):
            text = ax.text(j, i, int(table[i, j]),
                          ha="center", va="center", color="black",
                          fontweight='bold', fontsize=12)

    plt.colorbar(im1, ax=ax)

    # 2. Expected counts heatmap
    ax = axes[1]
    expected = result['expected_counts']
    im2 = ax.imshow(expected, cmap='Blues', aspect='auto')
    ax.set_xticks(range(expected.shape[1]))
    ax.set_yticks(range(expected.shape[0]))
    ax.set_xticklabels(col_labels)
    ax.set_yticklabels(row_labels)
    ax.set_title('Expected Counts (under H0)')

    # Add text annotations
    for i in range(expected.shape[0]):
        for j in range(expected.shape[1]):
            text = ax.text(j, i, f'{expected[i, j]:.1f}',
                          ha="center", va="center", color="black",
                          fontweight='bold', fontsize=12)

    plt.colorbar(im2, ax=ax)

    # 3. Standardized residuals
    ax = axes[2]
    residuals = (table - expected) / np.sqrt(expected)
    max_abs = np.max(np.abs(residuals))
    im3 = ax.imshow(residuals, cmap='RdBu_r', aspect='auto',
                    vmin=-max_abs, vmax=max_abs)
    ax.set_xticks(range(residuals.shape[1]))
    ax.set_yticks(range(residuals.shape[0]))
    ax.set_xticklabels(col_labels)
    ax.set_yticklabels(row_labels)
    ax.set_title('Standardized Residuals\n(|residual| > 2 suggests departure)')

    # Add text annotations
    for i in range(residuals.shape[0]):
        for j in range(residuals.shape[1]):
            color = 'white' if abs(residuals[i, j]) > 1.5 else 'black'
            text = ax.text(j, i, f'{residuals[i, j]:.2f}',
                          ha="center", va="center", color=color,
                          fontweight='bold', fontsize=10)

    plt.colorbar(im3, ax=ax)

    plt.suptitle(f'Chi-Squared Test: $\chi^2$ = {result["chi2_statistic"]:.2f}, '
                f'p = {result["p_value"]:.4f}',
                fontsize=14, fontweight='bold')
    plt.tight_layout()
    return fig

# Example usage
print("=== Chi-Squared Test Example ===\n")

# Example 1: Standard 2×2 table
table1 = np.array([[160, 1840],
                   [190, 1810]])

chi2_1, fisher_1 = compare_chi_square_fisher(table1, alpha=0.05)

# Verify equivalence with z-test
x1, n1 = table1[0, 0], table1[0].sum()
x2, n2 = table1[1, 0], table1[1].sum()

p1 = x1 / n1
p2 = x2 / n2
p_pooled = (x1 + x2) / (n1 + n2)

z = (p2 - p1) / np.sqrt(p_pooled * (1 - p_pooled) * (1/n1 + 1/n2))
p_z = 2 * (1 - stats.norm.cdf(abs(z)))

print("=== Equivalence Verification (2×2 table) ===")
print(f"Z-test p-value:        {p_z:.4f}")
print(f"Chi-squared p-value:   {chi2_1['p_value']:.4f}")
print(f"z^2 = {z**2:.3f}, chi^2 = {chi2_1['chi2_statistic']:.3f}")
print(f"Difference: {abs(z**2 - chi2_1['chi2_statistic']):.6f}\n")

# Example 2: Small expected counts (Fisher's better)
table2 = np.array([[8, 12],
                   [2, 18]])

print("=== Example 2: Small Expected Counts ===\n")
chi2_2, fisher_2 = compare_chi_square_fisher(table2, alpha=0.05)

# Example 3: Larger table (3×3)
table3 = np.array([[30, 20, 10],
                   [25, 30, 15],
                   [15, 25, 30]])

print("\n=== Example 3: 3×3 Table ===\n")
chi2_3, _ = compare_chi_square_fisher(table3, alpha=0.05)

# Visualize
row_labels = ['Control', 'Treatment']
col_labels = ['Converted', 'Not Converted']

fig1 = visualize_contingency_table(table1, chi2_1, row_labels, col_labels)
plt.savefig('chi_square_2x2.pdf')

category_labels = ['Category A', 'Category B', 'Category C']
fig3 = visualize_contingency_table(table3, chi2_3,
                                    category_labels, category_labels)
plt.savefig('chi_square_3x3.pdf')

plt.show()
\end{lstlisting}

\section{Understanding P-Values}

\begin{definition}[P-Value]
The p-value is the probability, under the null hypothesis, of observing a test statistic at least as extreme as the one actually observed.

Formally: $p = P(T \geq t_{\text{obs}} | H_0)$ for one-sided, or $p = P(|T| \geq |t_{\text{obs}}| | H_0)$ for two-sided.
\end{definition}

\subsection{Correct Interpretation}

A p-value of 0.03 means:
\begin{itemize}
    \item \checkmark If $H_0$ were true, we would observe data this extreme in 3\% of experiments
    \item \checkmark The data provide evidence against $H_0$ (but don't prove $H_1$)
    \item \checkmark At $\alpha = 0.05$, we reject $H_0$
\end{itemize}

\subsection{Common Misinterpretations}

\begin{enumerate}
    \item \textbf{WRONG:} "The probability that $H_0$ is true is 0.03"

    P-values are \textit{not} probabilities of hypotheses. They assume $H_0$ is true and calculate probability of data.

    Correct: $P(\text{data}|H_0)$, not $P(H_0|\text{data})$

    \item \textbf{WRONG:} "The probability that the result occurred by chance is 0.03"

    Everything has some probability. This confuses conditional probabilities and event probabilities.

    \item \textbf{WRONG:} "A p-value of 0.049 is meaningfully different from 0.051"

    The $\alpha = 0.05$ threshold is arbitrary. P-values are continuous measures of evidence against $H_0$, not binary outcomes.

    \item \textbf{WRONG:} "A non-significant result means no effect exists"

    Absence of evidence is not evidence of absence. Failure to reject $H_0$ may indicate:
    \begin{itemize}
        \item True null effect
        \item Insufficient power to detect existing effect
        \item High variance masking signal
    \end{itemize}

    \item \textbf{WRONG:} "A significant result means the effect is large or important"

    Statistical significance depends on sample size. With n=1,000,000, tiny meaningless effects become "significant."
\end{enumerate}

\subsection{Statistical Significance vs. Practical Significance}

Statistical significance ($p < \alpha$) does not imply practical importance.

\begin{example}[Large Sample Paradox]
With $n = 100{,}000$ per variant:

Conversion rate: 10.00\% → 10.05\% (0.05 percentage points)

This is highly statistically significant ($p < 0.001$) but may be economically negligible if implementation cost exceeds benefit.

Conversely, with $n = 100$:

Conversion rate: 10\% → 15\% (5 percentage points)

This may not reach significance ($p = 0.12$) despite clear practical importance.
\end{example}

\textbf{Recommendation:} Report effect sizes, confidence intervals, and business impact alongside p-values. Make decisions based on practical significance informed by statistical uncertainty.

\section{Assumptions and Diagnostics}

All parametric tests make assumptions. Violating assumptions can invalidate conclusions.

\subsection{Normality Testing}

Many tests (t-test, ANOVA) assume data are approximately normally distributed.

\textbf{Visual diagnostics:}
\begin{itemize}
    \item Histograms
    \item Q-Q plots (quantile-quantile)
    \item Kernel density estimates
\end{itemize}

\textbf{Formal tests:}
\begin{itemize}
    \item Shapiro-Wilk test (most powerful for small-medium samples)
    \item Kolmogorov-Smirnov test
    \item Anderson-Darling test
\end{itemize}

\textbf{Important caveat:} With large samples, normality tests often reject $H_0$ for trivial departures. Visual inspection often more useful.

\textbf{Robustness:} T-tests are robust to moderate departures from normality, especially with:
\begin{itemize}
    \item Large samples ($n > 30$ per group)
    \item Symmetric distributions
    \item Equal sample sizes
\end{itemize}

\subsection{Variance Homogeneity}

Pooled t-test and ANOVA assume equal variances across groups.

\textbf{Visual diagnostics:}
\begin{itemize}
    \item Side-by-side boxplots
    \item Residual plots
    \item Spread-location plots
\end{itemize}

\textbf{Formal tests:}
\begin{itemize}
    \item Levene's test (robust to non-normality)
    \item Bartlett's test (sensitive to non-normality)
    \item F-test (assumes normality)
\end{itemize}

\textbf{Rule of thumb:} If largest variance / smallest variance $> 4$, consider variances unequal.

\textbf{Solution:} Use Welch's t-test (doesn't assume equal variances).

\subsection{Independence Assumptions}

Most critical assumption, hardest to test.

\textbf{Violations:}
\begin{itemize}
    \item Repeated measures (same user appears multiple times)
    \item Clustering (users within groups more similar)
    \item Time series (temporal autocorrelation)
    \item Spatial correlation
    \item Network effects
\end{itemize}

\textbf{Detection:}
\begin{itemize}
    \item Check data collection process
    \item Autocorrelation plots (time series)
    \item Intraclass correlation coefficient (clustering)
\end{itemize}

\textbf{Solutions:}
\begin{itemize}
    \item Mixed effects models
    \item Cluster-robust standard errors
    \item Time series methods
    \item Aggregation to independent units
\end{itemize}

\subsection{Outlier Detection and Treatment}

Outliers can severely affect means and standard deviations, distorting test results.

\textbf{Detection methods:}
\begin{itemize}
    \item Boxplot (IQR method): values $< Q_1 - 1.5 \times IQR$ or $> Q_3 + 1.5 \times IQR$
    \item Z-score: $|z| > 3$ (or 2.5 for stricter criterion)
    \item Modified Z-score using median absolute deviation (robust)
    \item Visual inspection (always!)
\end{itemize}

\textbf{Treatment options:}
\begin{enumerate}
    \item \textbf{Investigate:} Are outliers errors or valid extreme values?

    \item \textbf{Cap/winsorize:} Replace extreme values with threshold (e.g., 99th percentile)

    \item \textbf{Transform:} Log transform to reduce impact of extremes

    \item \textbf{Remove:} Only if clear data errors, document clearly

    \item \textbf{Robust methods:} Use median instead of mean, trimmed means, non-parametric tests

    \item \textbf{Separate analysis:} Analyze with and without outliers, report both
\end{enumerate}

\textbf{Never:} Selectively remove outliers to achieve desired p-value.

\lstset{language=Python, caption=Comprehensive Diagnostic Functions}
\begin{lstlisting}
import numpy as np
from scipy import stats
import matplotlib.pyplot as plt
import seaborn as sns
from typing import Dict, Tuple
import pandas as pd

def check_normality(
    data: np.ndarray,
    name: str = "Data",
    alpha: float = 0.05
) -> Dict:
    """
    Check normality assumption using multiple methods.

    Parameters:
    -----------
    data : np.ndarray
        Data to test
    name : str
        Name for labeling
    alpha : float
        Significance level for tests

    Returns:
    --------
    dict : Normality test results
    """
    # Shapiro-Wilk test
    shapiro_stat, shapiro_p = stats.shapiro(data)

    # Kolmogorov-Smirnov test
    ks_stat, ks_p = stats.kstest(data, 'norm',
                                  args=(np.mean(data), np.std(data)))

    # Anderson-Darling test
    anderson_result = stats.anderson(data, dist='norm')

    # Skewness and kurtosis
    skewness = stats.skew(data)
    kurtosis = stats.kurtosis(data)

    return {
        'name': name,
        'n': len(data),
        'shapiro_statistic': shapiro_stat,
        'shapiro_pvalue': shapiro_p,
        'shapiro_normal': shapiro_p > alpha,
        'ks_statistic': ks_stat,
        'ks_pvalue': ks_p,
        'ks_normal': ks_p > alpha,
        'anderson_statistic': anderson_result.statistic,
        'anderson_critical_values': anderson_result.critical_values,
        'skewness': skewness,
        'kurtosis': kurtosis
    }

def check_variance_equality(
    data1: np.ndarray,
    data2: np.ndarray,
    alpha: float = 0.05
) -> Dict:
    """
    Check equality of variances.

    Parameters:
    -----------
    data1, data2 : np.ndarray
        Data for each group
    alpha : float
        Significance level

    Returns:
    --------
    dict : Variance equality test results
    """
    var1, var2 = np.var(data1, ddof=1), np.var(data2, ddof=1)
    std1, std2 = np.std(data1, ddof=1), np.std(data2, ddof=1)

    # Variance ratio
    var_ratio = max(var1, var2) / min(var1, var2)

    # Levene's test (robust to non-normality)
    levene_stat, levene_p = stats.levene(data1, data2)

    # Bartlett's test (assumes normality)
    bartlett_stat, bartlett_p = stats.bartlett(data1, data2)

    # F-test (assumes normality)
    if var1 >= var2:
        f_stat = var1 / var2
        df1, df2 = len(data1) - 1, len(data2) - 1
    else:
        f_stat = var2 / var1
        df2, df1 = len(data1) - 1, len(data2) - 1

    f_p = 2 * min(stats.f.cdf(f_stat, df1, df2),
                  1 - stats.f.cdf(f_stat, df1, df2))

    return {
        'var1': var1,
        'var2': var2,
        'std1': std1,
        'std2': std2,
        'variance_ratio': var_ratio,
        'levene_statistic': levene_stat,
        'levene_pvalue': levene_p,
        'levene_equal': levene_p > alpha,
        'bartlett_statistic': bartlett_stat,
        'bartlett_pvalue': bartlett_p,
        'bartlett_equal': bartlett_p > alpha,
        'f_statistic': f_stat,
        'f_pvalue': f_p,
        'f_equal': f_p > alpha
    }

def detect_outliers(
    data: np.ndarray,
    method: str = 'iqr'
) -> Dict:
    """
    Detect outliers using various methods.

    Parameters:
    -----------
    data : np.ndarray
        Data to analyze
    method : str
        'iqr', 'zscore', or 'modified_zscore'

    Returns:
    --------
    dict : Outlier detection results
    """
    if method == 'iqr':
        q1 = np.percentile(data, 25)
        q3 = np.percentile(data, 75)
        iqr = q3 - q1
        lower_bound = q1 - 1.5 * iqr
        upper_bound = q3 + 1.5 * iqr
        outliers = (data < lower_bound) | (data > upper_bound)

    elif method == 'zscore':
        z_scores = np.abs(stats.zscore(data))
        outliers = z_scores > 3

    elif method == 'modified_zscore':
        median = np.median(data)
        mad = np.median(np.abs(data - median))
        modified_z_scores = 0.6745 * (data - median) / mad
        outliers = np.abs(modified_z_scores) > 3.5

    else:
        raise ValueError("method must be 'iqr', 'zscore', or 'modified_zscore'")

    outlier_indices = np.where(outliers)[0]
    outlier_values = data[outliers]

    return {
        'method': method,
        'n_outliers': np.sum(outliers),
        'outlier_proportion': np.mean(outliers),
        'outlier_indices': outlier_indices,
        'outlier_values': outlier_values,
        'outlier_mask': outliers
    }

def visualize_diagnostics(
    data1: np.ndarray,
    data2: np.ndarray,
    name1: str = "Control",
    name2: str = "Treatment"
):
    """Create comprehensive diagnostic plots."""
    fig = plt.figure(figsize=(16, 12))
    gs = fig.add_gridspec(3, 3, hspace=0.3, wspace=0.3)

    # 1. Histograms with normal overlay
    ax1 = fig.add_subplot(gs[0, 0])
    ax1.hist(data1, bins=30, alpha=0.6, density=True, label=name1, color='C0')
    mu1, std1 = np.mean(data1), np.std(data1)
    x1 = np.linspace(data1.min(), data1.max(), 100)
    ax1.plot(x1, stats.norm.pdf(x1, mu1, std1), 'C0--', linewidth=2,
            label='Normal fit')
    ax1.set_xlabel('Value')
    ax1.set_ylabel('Density')
    ax1.set_title(f'{name1} Distribution')
    ax1.legend()
    ax1.grid(alpha=0.3)

    ax2 = fig.add_subplot(gs[0, 1])
    ax2.hist(data2, bins=30, alpha=0.6, density=True, label=name2, color='C1')
    mu2, std2 = np.mean(data2), np.std(data2)
    x2 = np.linspace(data2.min(), data2.max(), 100)
    ax2.plot(x2, stats.norm.pdf(x2, mu2, std2), 'C1--', linewidth=2,
            label='Normal fit')
    ax2.set_xlabel('Value')
    ax2.set_ylabel('Density')
    ax2.set_title(f'{name2} Distribution')
    ax2.legend()
    ax2.grid(alpha=0.3)

    # 2. Q-Q plots
    ax3 = fig.add_subplot(gs[0, 2])
    stats.probplot(data1, dist="norm", plot=ax3)
    ax3.set_title(f'{name1} Q-Q Plot')
    ax3.grid(alpha=0.3)

    ax4 = fig.add_subplot(gs[1, 0])
    stats.probplot(data2, dist="norm", plot=ax4)
    ax4.set_title(f'{name2} Q-Q Plot')
    ax4.grid(alpha=0.3)

    # 3. Box plots
    ax5 = fig.add_subplot(gs[1, 1])
    box_data = [data1, data2]
    bp = ax5.boxplot(box_data, labels=[name1, name2],
                     patch_artist=True, widths=0.6)
    for patch, color in zip(bp['boxes'], ['C0', 'C1']):
        patch.set_facecolor(color)
        patch.set_alpha(0.6)
    ax5.set_ylabel('Value')
    ax5.set_title('Box Plots (outlier detection)')
    ax5.grid(axis='y', alpha=0.3)

    # Add variance ratio annotation
    var_ratio = max(np.var(data1), np.var(data2)) / min(np.var(data1), np.var(data2))
    ax5.text(0.5, 0.95, f'Variance ratio: {var_ratio:.2f}',
            transform=ax5.transAxes, ha='center', va='top',
            bbox=dict(boxstyle='round', facecolor='wheat', alpha=0.5))

    # 4. Violin plots with individual points
    ax6 = fig.add_subplot(gs[1, 2])
    parts = ax6.violinplot([data1, data2], positions=[1, 2],
                           showmeans=True, showmedians=True)
    ax6.set_xticks([1, 2])
    ax6.set_xticklabels([name1, name2])
    ax6.set_ylabel('Value')
    ax6.set_title('Violin Plots')
    ax6.grid(axis='y', alpha=0.3)

    # 5. ECDF comparison
    ax7 = fig.add_subplot(gs[2, 0])
    sorted_data1 = np.sort(data1)
    sorted_data2 = np.sort(data2)
    ecdf1 = np.arange(1, len(sorted_data1) + 1) / len(sorted_data1)
    ecdf2 = np.arange(1, len(sorted_data2) + 1) / len(sorted_data2)
    ax7.plot(sorted_data1, ecdf1, label=name1, linewidth=2, color='C0')
    ax7.plot(sorted_data2, ecdf2, label=name2, linewidth=2, color='C1')
    ax7.set_xlabel('Value')
    ax7.set_ylabel('ECDF')
    ax7.set_title('Empirical Cumulative Distribution')
    ax7.legend()
    ax7.grid(alpha=0.3)

    # 6. Variance comparison
    ax8 = fig.add_subplot(gs[2, 1])
    variances = [np.var(data1, ddof=1), np.var(data2, ddof=1)]
    std_devs = [np.std(data1, ddof=1), np.std(data2, ddof=1)]

    x_pos = np.arange(2)
    bars = ax8.bar(x_pos, std_devs, alpha=0.7, color=['C0', 'C1'])
    ax8.set_xticks(x_pos)
    ax8.set_xticklabels([name1, name2])
    ax8.set_ylabel('Standard Deviation')
    ax8.set_title('Standard Deviations Comparison')
    ax8.grid(axis='y', alpha=0.3)

    for i, (bar, std) in enumerate(zip(bars, std_devs)):
        height = bar.get_height()
        ax8.text(bar.get_x() + bar.get_width()/2., height,
                f'{std:.2f}', ha='center', va='bottom', fontweight='bold')

    # 7. Summary statistics table
    ax9 = fig.add_subplot(gs[2, 2])
    ax9.axis('off')

    summary_data = [
        ['Statistic', name1, name2],
        ['n', f'{len(data1)}', f'{len(data2)}'],
        ['Mean', f'{np.mean(data1):.2f}', f'{np.mean(data2):.2f}'],
        ['Median', f'{np.median(data1):.2f}', f'{np.median(data2):.2f}'],
        ['Std Dev', f'{np.std(data1, ddof=1):.2f}', f'{np.std(data2, ddof=1):.2f}'],
        ['Min', f'{data1.min():.2f}', f'{data2.min():.2f}'],
        ['Max', f'{data1.max():.2f}', f'{data2.max():.2f}'],
        ['Q1', f'{np.percentile(data1, 25):.2f}', f'{np.percentile(data2, 25):.2f}'],
        ['Q3', f'{np.percentile(data1, 75):.2f}', f'{np.percentile(data2, 75):.2f}'],
        ['Skewness', f'{stats.skew(data1):.2f}', f'{stats.skew(data2):.2f}'],
        ['Kurtosis', f'{stats.kurtosis(data1):.2f}', f'{stats.kurtosis(data2):.2f}']
    ]

    table = ax9.table(cellText=summary_data, cellLoc='center',
                     loc='center', bbox=[0, 0, 1, 1])
    table.auto_set_font_size(False)
    table.set_fontsize(9)
    table.scale(1, 1.8)

    # Style header row
    for i in range(3):
        table[(0, i)].set_facecolor('#4CAF50')
        table[(0, i)].set_text_props(weight='bold', color='white')

    plt.suptitle('Comprehensive Diagnostic Plots', fontsize=16, fontweight='bold')

    return fig

# Example usage
print("=== Assumption Diagnostics ===\n")

np.random.seed(42)

# Generate test data
data1_normal = np.random.normal(50, 10, 500)
data2_normal = np.random.normal(53, 10, 500)

data1_skewed = np.random.gamma(2, 5, 500)
data2_skewed = np.random.gamma(2.5, 5, 500)

data1_unequal_var = np.random.normal(50, 10, 500)
data2_unequal_var = np.random.normal(53, 20, 500)

# Check normality
print("=== Normality Tests (Normal Data) ===")
norm_result1 = check_normality(data1_normal, "Control (Normal)")
print(f"{norm_result1['name']} (n={norm_result1['n']}):")
print(f"  Shapiro-Wilk: W={norm_result1['shapiro_statistic']:.4f}, "
      f"p={norm_result1['shapiro_pvalue']:.4f}, "
      f"Normal: {norm_result1['shapiro_normal']}")
print(f"  Skewness: {norm_result1['skewness']:.3f}")
print(f"  Kurtosis: {norm_result1['kurtosis']:.3f}\n")

print("=== Normality Tests (Skewed Data) ===")
norm_result2 = check_normality(data1_skewed, "Control (Skewed)")
print(f"{norm_result2['name']} (n={norm_result2['n']}):")
print(f"  Shapiro-Wilk: W={norm_result2['shapiro_statistic']:.4f}, "
      f"p={norm_result2['shapiro_pvalue']:.4e}, "
      f"Normal: {norm_result2['shapiro_normal']}")
print(f"  Skewness: {norm_result2['skewness']:.3f}")
print(f"  Kurtosis: {norm_result2['kurtosis']:.3f}\n")

# Check variance equality
print("=== Variance Equality Tests ===")
var_result = check_variance_equality(data1_unequal_var, data2_unequal_var)
print(f"Variance 1: {var_result['var1']:.2f}, Variance 2: {var_result['var2']:.2f}")
print(f"Variance ratio: {var_result['variance_ratio']:.2f}")
print(f"Levene's test: F={var_result['levene_statistic']:.4f}, "
      f"p={var_result['levene_pvalue']:.4f}, "
      f"Equal: {var_result['levene_equal']}\n")

# Detect outliers
print("=== Outlier Detection ===")
# Add some outliers
data_with_outliers = np.concatenate([data1_normal, [100, 105, -20]])
outlier_result = detect_outliers(data_with_outliers, method='iqr')
print(f"Method: {outlier_result['method']}")
print(f"Number of outliers: {outlier_result['n_outliers']}")
print(f"Outlier proportion: {outlier_result['outlier_proportion']:.3f}")
print(f"Outlier values: {outlier_result['outlier_values']}\n")

# Visualize
fig1 = visualize_diagnostics(data1_normal, data2_normal)
plt.savefig('diagnostics_normal.pdf')

fig2 = visualize_diagnostics(data1_skewed, data2_skewed,
                             "Control (Skewed)", "Treatment (Skewed)")
plt.savefig('diagnostics_skewed.pdf')

fig3 = visualize_diagnostics(data1_unequal_var, data2_unequal_var,
                             "Control (σ=10)", "Treatment (σ=20)")
plt.savefig('diagnostics_unequal_var.pdf')

plt.show()
\end{lstlisting}

\section{Non-Parametric Alternatives}

When parametric assumptions fail, non-parametric tests provide robust alternatives.

\subsection{Mann-Whitney U Test}

The Mann-Whitney U test (also called Wilcoxon rank-sum test) tests whether two samples come from populations with the same distribution.

\textbf{Null hypothesis:} The distributions of the two groups are identical (specifically, $P(X_1 > X_2) = 0.5$).

\textbf{Procedure:}
\begin{enumerate}
    \item Rank all observations from both groups combined
    \item Calculate sum of ranks for each group: $R_1$, $R_2$
    \item Calculate U statistic:
    \begin{equation}
    U_1 = R_1 - \frac{n_1(n_1 + 1)}{2}, \quad U_2 = R_2 - \frac{n_2(n_2 + 1)}{2}
    \end{equation}
    \item Use $U = \min(U_1, U_2)$ as test statistic
\end{enumerate}

\textbf{Advantages:}
\begin{itemize}
    \item No normality assumption
    \item Robust to outliers
    \item Works for ordinal data
    \item More powerful than t-test for heavy-tailed distributions
\end{itemize}

\textbf{Disadvantages:}
\begin{itemize}
    \item Less powerful than t-test when data are normal
    \item Tests distribution similarity, not just location
    \item Harder to interpret (medians, not means)
\end{itemize}

\subsection{Permutation Tests}

Permutation tests provide exact p-values without distributional assumptions.

\textbf{Principle:} Under $H_0$ of no treatment effect, treatment labels are exchangeable.

\textbf{Procedure:}
\begin{enumerate}
    \item Calculate observed test statistic $T_{\text{obs}}$
    \item Randomly permute treatment labels many times (e.g., 10,000)
    \item For each permutation, calculate test statistic
    \item P-value = proportion of permutations with $|T| \geq |T_{\text{obs}}|$
\end{enumerate}

\textbf{Advantages:}
\begin{itemize}
    \item Exact finite-sample validity (no approximations)
    \item No distributional assumptions
    \item Works for any test statistic (means, medians, quantiles, etc.)
    \item Intuitive interpretation
\end{itemize}

\textbf{Disadvantages:}
\begin{itemize}
    \item Computationally intensive
    \item Requires exchangeability (randomization)
    \item Less familiar to stakeholders
\end{itemize}

\subsection{Bootstrap Methods}

Bootstrap resampling estimates sampling distributions and constructs confidence intervals.

\textbf{Procedure for confidence interval:}
\begin{enumerate}
    \item Resample with replacement from each group
    \item Calculate statistic of interest (e.g., difference in means)
    \item Repeat B times (e.g., B=10,000)
    \item Use percentiles of bootstrap distribution for CI
\end{enumerate}

\textbf{Percentile CI:}
\begin{equation}
\text{CI}_{1-\alpha} = [\text{Percentile}_{\alpha/2}, \text{Percentile}_{1-\alpha/2}]
\end{equation}

\textbf{Bias-corrected and accelerated (BCa) CI:} Adjusts for bias and skewness in bootstrap distribution.

\lstset{language=Python, caption=Non-Parametric Tests Implementation}
\begin{lstlisting}
import numpy as np
from scipy import stats
import matplotlib.pyplot as plt
from typing import Dict
import pandas as pd

def mann_whitney_test(
    data1: np.ndarray,
    data2: np.ndarray,
    alpha: float = 0.05,
    alternative: str = 'two-sided'
) -> Dict:
    """
    Perform Mann-Whitney U test.

    Parameters:
    -----------
    data1, data2 : np.ndarray
        Data for each group
    alpha : float
        Significance level
    alternative : str
        'two-sided', 'greater', or 'less'

    Returns:
    --------
    dict : Test results
    """
    u_stat, p_value = stats.mannwhitneyu(
        data1, data2, alternative=alternative
    )

    reject_null = p_value < alpha

    # Effect size: rank-biserial correlation
    n1, n2 = len(data1), len(data2)
    rank_biserial = 1 - (2*u_stat) / (n1 * n2)

    # Median difference (approximate)
    median1, median2 = np.median(data1), np.median(data2)
    median_diff = median2 - median1

    return {
        'test': 'Mann-Whitney U test',
        'u_statistic': u_stat,
        'p_value': p_value,
        'alpha': alpha,
        'reject_null': reject_null,
        'alternative': alternative,
        'median1': median1,
        'median2': median2,
        'median_diff': median_diff,
        'rank_biserial': rank_biserial,
        'n1': n1,
        'n2': n2,
        'decision': 'Reject H0' if reject_null else 'Fail to reject H0'
    }

def permutation_test(
    data1: np.ndarray,
    data2: np.ndarray,
    n_permutations: int = 10000,
    statistic: str = 'mean_diff',
    alternative: str = 'two-sided',
    seed: int = None
) -> Dict:
    """
    Perform permutation test.

    Parameters:
    -----------
    data1, data2 : np.ndarray
        Data for each group
    n_permutations : int
        Number of permutations
    statistic : str
        Test statistic: 'mean_diff', 'median_diff', 't_stat', or custom function
    alternative : str
        'two-sided', 'greater', or 'less'
    seed : int
        Random seed

    Returns:
    --------
    dict : Test results
    """
    if seed is not None:
        np.random.seed(seed)

    combined = np.concatenate([data1, data2])
    n1 = len(data1)
    n_total = len(combined)

    # Define test statistic function
    def calc_statistic(d1, d2):
        if statistic == 'mean_diff':
            return np.mean(d2) - np.mean(d1)
        elif statistic == 'median_diff':
            return np.median(d2) - np.median(d1)
        elif statistic == 't_stat':
            # Welch's t-statistic
            m1, m2 = np.mean(d1), np.mean(d2)
            v1, v2 = np.var(d1, ddof=1), np.var(d2, ddof=1)
            n1, n2 = len(d1), len(d2)
            return (m2 - m1) / np.sqrt(v1/n1 + v2/n2)
        elif callable(statistic):
            return statistic(d1, d2)
        else:
            raise ValueError("Invalid statistic")

    # Observed test statistic
    obs_stat = calc_statistic(data1, data2)

    # Permutation distribution
    perm_stats = np.zeros(n_permutations)

    for i in range(n_permutations):
        # Randomly permute
        permuted = np.random.permutation(combined)
        perm_data1 = permuted[:n1]
        perm_data2 = permuted[n1:]

        perm_stats[i] = calc_statistic(perm_data1, perm_data2)

    # P-value
    if alternative == 'two-sided':
        p_value = np.mean(np.abs(perm_stats) >= np.abs(obs_stat))
    elif alternative == 'greater':
        p_value = np.mean(perm_stats >= obs_stat)
    elif alternative == 'less':
        p_value = np.mean(perm_stats <= obs_stat)
    else:
        raise ValueError("alternative must be 'two-sided', 'greater', or 'less'")

    reject_null = p_value < 0.05

    return {
        'test': 'Permutation test',
        'observed_statistic': obs_stat,
        'permutation_statistics': perm_stats,
        'p_value': p_value,
        'n_permutations': n_permutations,
        'alternative': alternative,
        'reject_null': reject_null,
        'statistic_type': statistic if isinstance(statistic, str) else 'custom',
        'decision': 'Reject H0' if reject_null else 'Fail to reject H0'
    }

def bootstrap_ci(
    data1: np.ndarray,
    data2: np.ndarray,
    statistic: str = 'mean_diff',
    n_bootstrap: int = 10000,
    alpha: float = 0.05,
    method: str = 'percentile',
    seed: int = None
) -> Dict:
    """
    Bootstrap confidence interval for difference.

    Parameters:
    -----------
    data1, data2 : np.ndarray
        Data for each group
    statistic : str
        'mean_diff', 'median_diff', or custom function
    n_bootstrap : int
        Number of bootstrap samples
    alpha : float
        Significance level
    method : str
        'percentile' or 'bca' (bias-corrected accelerated)
    seed : int
        Random seed

    Returns:
    --------
    dict : Bootstrap results
    """
    if seed is not None:
        np.random.seed(seed)

    def calc_statistic(d1, d2):
        if statistic == 'mean_diff':
            return np.mean(d2) - np.mean(d1)
        elif statistic == 'median_diff':
            return np.median(d2) - np.median(d1)
        elif callable(statistic):
            return statistic(d1, d2)
        else:
            raise ValueError("Invalid statistic")

    # Observed statistic
    obs_stat = calc_statistic(data1, data2)

    # Bootstrap distribution
    bootstrap_stats = np.zeros(n_bootstrap)

    for i in range(n_bootstrap):
        # Resample with replacement from each group
        boot_data1 = np.random.choice(data1, size=len(data1), replace=True)
        boot_data2 = np.random.choice(data2, size=len(data2), replace=True)

        bootstrap_stats[i] = calc_statistic(boot_data1, boot_data2)

    if method == 'percentile':
        ci_lower = np.percentile(bootstrap_stats, 100 * alpha/2)
        ci_upper = np.percentile(bootstrap_stats, 100 * (1 - alpha/2))

    elif method == 'bca':
        # Bias correction
        z0 = stats.norm.ppf(np.mean(bootstrap_stats < obs_stat))

        # Acceleration (jackknife)
        n1, n2 = len(data1), len(data2)
        jack_stats = []

        for i in range(n1):
            jack_data1 = np.delete(data1, i)
            jack_stats.append(calc_statistic(jack_data1, data2))

        jack_mean = np.mean(jack_stats)
        numerator = np.sum((jack_mean - jack_stats)**3)
        denominator = 6 * (np.sum((jack_mean - jack_stats)**2)**(3/2))
        a = numerator / denominator if denominator != 0 else 0

        # Adjusted percentiles
        z_alpha_2 = stats.norm.ppf(alpha/2)
        z_1_alpha_2 = stats.norm.ppf(1 - alpha/2)

        alpha_1 = stats.norm.cdf(z0 + (z0 + z_alpha_2)/(1 - a*(z0 + z_alpha_2)))
        alpha_2 = stats.norm.cdf(z0 + (z0 + z_1_alpha_2)/(1 - a*(z0 + z_1_alpha_2)))

        ci_lower = np.percentile(bootstrap_stats, 100 * alpha_1)
        ci_upper = np.percentile(bootstrap_stats, 100 * alpha_2)

    else:
        raise ValueError("method must be 'percentile' or 'bca'")

    # Standard error
    bootstrap_se = np.std(bootstrap_stats)

    return {
        'method': f'Bootstrap {method} CI',
        'observed_statistic': obs_stat,
        'bootstrap_statistics': bootstrap_stats,
        'bootstrap_se': bootstrap_se,
        'ci_lower': ci_lower,
        'ci_upper': ci_upper,
        'alpha': alpha,
        'n_bootstrap': n_bootstrap,
        'statistic_type': statistic if isinstance(statistic, str) else 'custom'
    }

def visualize_nonparametric_tests(
    data1: np.ndarray,
    data2: np.ndarray,
    perm_result: Dict,
    boot_result: Dict,
    name1: str = "Control",
    name2: str = "Treatment"
):
    """Visualize non-parametric test results."""
    fig, axes = plt.subplots(2, 2, figsize=(14, 10))

    # 1. Permutation distribution
    ax = axes[0, 0]
    ax.hist(perm_result['permutation_statistics'], bins=50, density=True,
           alpha=0.7, label='Permutation distribution', color='skyblue')
    ax.axvline(perm_result['observed_statistic'], color='red',
              linestyle='--', linewidth=2,
              label=f'Observed = {perm_result["observed_statistic"]:.3f}')

    # Add p-value region
    if perm_result['alternative'] == 'two-sided':
        obs_abs = abs(perm_result['observed_statistic'])
        ax.axvline(-obs_abs, color='red', linestyle='--',
                  linewidth=2, alpha=0.5)

    ax.set_xlabel('Test Statistic')
    ax.set_ylabel('Density')
    ax.set_title(f'Permutation Test\n'
                f'p-value = {perm_result["p_value"]:.4f}')
    ax.legend()
    ax.grid(alpha=0.3)

    # 2. Bootstrap distribution
    ax = axes[0, 1]
    ax.hist(boot_result['bootstrap_statistics'], bins=50, density=True,
           alpha=0.7, label='Bootstrap distribution', color='lightgreen')
    ax.axvline(boot_result['observed_statistic'], color='red',
              linestyle='--', linewidth=2,
              label=f'Observed = {boot_result["observed_statistic"]:.3f}')
    ax.axvline(boot_result['ci_lower'], color='orange',
              linestyle=':', linewidth=2,
              label=f'CI: [{boot_result["ci_lower"]:.3f}, {boot_result["ci_upper"]:.3f}]')
    ax.axvline(boot_result['ci_upper'], color='orange',
              linestyle=':', linewidth=2)

    ax.set_xlabel('Difference')
    ax.set_ylabel('Density')
    ax.set_title(f'Bootstrap Distribution\n'
                f'SE = {boot_result["bootstrap_se"]:.3f}')
    ax.legend()
    ax.grid(alpha=0.3)

    # 3. ECDF comparison (Mann-Whitney concept)
    ax = axes[1, 0]
    sorted_data1 = np.sort(data1)
    sorted_data2 = np.sort(data2)
    ecdf1 = np.arange(1, len(sorted_data1) + 1) / len(sorted_data1)
    ecdf2 = np.arange(1, len(sorted_data2) + 1) / len(sorted_data2)
    ax.plot(sorted_data1, ecdf1, label=name1, linewidth=2, color='C0')
    ax.plot(sorted_data2, ecdf2, label=name2, linewidth=2, color='C1')

    # Add medians
    ax.axvline(np.median(data1), color='C0', linestyle='--', alpha=0.5)
    ax.axvline(np.median(data2), color='C1', linestyle='--', alpha=0.5)

    ax.set_xlabel('Value')
    ax.set_ylabel('ECDF')
    ax.set_title('Empirical CDFs (Mann-Whitney Concept)')
    ax.legend()
    ax.grid(alpha=0.3)

    # 4. Comparison table
    ax = axes[1, 1]
    ax.axis('off')

    # Run all tests for comparison
    t_result = t_test_means(data1, data2)
    mw_result = mann_whitney_test(data1, data2)

    comparison_data = [
        ['Test', 'Statistic', 'P-value', 'Decision'],
        ['T-test', f'{t_result["t_statistic"]:.3f}',
         f'{t_result["p_value"]:.4f}', t_result['decision']],
        ['Mann-Whitney', f'{mw_result["u_statistic"]:.1f}',
         f'{mw_result["p_value"]:.4f}', mw_result['decision']],
        ['Permutation', f'{perm_result["observed_statistic"]:.3f}',
         f'{perm_result["p_value"]:.4f}', perm_result['decision']],
        ['Bootstrap CI', f'[{boot_result["ci_lower"]:.3f}, {boot_result["ci_upper"]:.3f}]',
         '--', 'CI includes 0' if boot_result['ci_lower'] < 0 < boot_result['ci_upper'] else 'CI excludes 0']
    ]

    table = ax.table(cellText=comparison_data, cellLoc='left',
                    loc='center', bbox=[0, 0, 1, 1])
    table.auto_set_font_size(False)
    table.set_fontsize(9)
    table.scale(1, 2.5)

    # Style header row
    for i in range(4):
        table[(0, i)].set_facecolor('#4CAF50')
        table[(0, i)].set_text_props(weight='bold', color='white')

    plt.suptitle('Non-Parametric Test Comparison', fontsize=14, fontweight='bold')
    plt.tight_layout()

    return fig

# Example usage
print("=== Non-Parametric Tests ===\n")

np.random.seed(42)

# Example 1: Normal data (t-test and non-parametric should agree)
data1_normal = np.random.normal(50, 10, 200)
data2_normal = np.random.normal(53, 10, 200)

# Example 2: Skewed data (non-parametric more appropriate)
data1_skewed = np.random.gamma(2, 5, 200)
data2_skewed = np.random.gamma(2.5, 5, 200)

print("=== Example 1: Normal Data ===\n")

# T-test
t_result = t_test_means(data1_normal, data2_normal)
print(f"T-test: t = {t_result['t_statistic']:.3f}, p = {t_result['p_value']:.4f}\n")

# Mann-Whitney
mw_result = mann_whitney_test(data1_normal, data2_normal)
print(f"Mann-Whitney: U = {mw_result['u_statistic']:.1f}, p = {mw_result['p_value']:.4f}")
print(f"Median difference: {mw_result['median_diff']:.3f}")
print(f"Rank-biserial correlation: {mw_result['rank_biserial']:.3f}\n")

# Permutation test
perm_result = permutation_test(data1_normal, data2_normal,
                               n_permutations=10000, seed=42)
print(f"Permutation test: p = {perm_result['p_value']:.4f}\n")

# Bootstrap CI
boot_result = bootstrap_ci(data1_normal, data2_normal,
                          n_bootstrap=10000, seed=42)
print(f"Bootstrap 95% CI: [{boot_result['ci_lower']:.3f}, {boot_result['ci_upper']:.3f}]")
print(f"Bootstrap SE: {boot_result['bootstrap_se']:.3f}\n")

print("=== Example 2: Skewed Data ===\n")

t_result_skewed = t_test_means(data1_skewed, data2_skewed)
mw_result_skewed = mann_whitney_test(data1_skewed, data2_skewed)
perm_result_skewed = permutation_test(data1_skewed, data2_skewed,
                                      n_permutations=10000, seed=42)
boot_result_skewed = bootstrap_ci(data1_skewed, data2_skewed,
                                  n_bootstrap=10000, seed=42)

print(f"T-test: t = {t_result_skewed['t_statistic']:.3f}, "
      f"p = {t_result_skewed['p_value']:.4f}")
print(f"Mann-Whitney: U = {mw_result_skewed['u_statistic']:.1f}, "
      f"p = {mw_result_skewed['p_value']:.4f}")
print(f"Permutation test: p = {perm_result_skewed['p_value']:.4f}")
print(f"Bootstrap 95% CI: [{boot_result_skewed['ci_lower']:.3f}, "
      f"{boot_result_skewed['ci_upper']:.3f}]\n")

# Visualize
fig1 = visualize_nonparametric_tests(data1_normal, data2_normal,
                                     perm_result, boot_result)
plt.savefig('nonparametric_normal.pdf')

fig2 = visualize_nonparametric_tests(data1_skewed, data2_skewed,
                                     perm_result_skewed, boot_result_skewed,
                                     "Control (Skewed)", "Treatment (Skewed)")
plt.savefig('nonparametric_skewed.pdf')

plt.show()
\end{lstlisting}

\subsection{When to Use Non-Parametric Tests}

\textbf{Use non-parametric tests when:}
\begin{itemize}
    \item Data are heavily skewed or have extreme outliers
    \item Sample sizes are small ($n < 30$ per group)
    \item Distributions are clearly non-normal (confirmed visually)
    \item Ordinal data (rankings, Likert scales)
    \item Variances are very unequal and sample sizes differ
\end{itemize}

\textbf{Use parametric tests when:}
\begin{itemize}
    \item Data are approximately normal
    \item Sample sizes are large ($n > 30$)
    \item More statistical power needed
    \item Easier interpretation desired (means vs. medians)
\end{itemize}

\textbf{Consider both:} For borderline cases, run both parametric and non-parametric tests. If conclusions agree, confidence increases. If they differ, investigate further.

\section{Test Selection Framework}

\lstset{language=Python, caption=Test Selection Decision Tree}
\begin{lstlisting}
def select_statistical_test(
    data1: np.ndarray,
    data2: np.ndarray,
    metric_type: str,
    check_assumptions: bool = True,
    alpha: float = 0.05
) -> Dict:
    """
    Recommend appropriate statistical test based on data characteristics.

    Parameters:
    -----------
    data1, data2 : np.ndarray
        Data for each group
    metric_type : str
        'proportion', 'mean', 'count', or 'ordinal'
    check_assumptions : bool
        Whether to check parametric assumptions
    alpha : float
        Significance level

    Returns:
    --------
    dict : Test recommendation and results
    """
    n1, n2 = len(data1), len(data2)

    recommendations = []
    tests_run = {}

    # Proportion test
    if metric_type == 'proportion':
        # Check sample size requirements
        if isinstance(data1[0], (int, np.integer)):
            # Binary data (0/1)
            x1, x2 = data1.sum(), data2.sum()
            p1, p2 = x1/n1, x2/n2

            min_count = min(x1, n1-x1, x2, n2-x2)

            if min_count >= 10:
                # Z-test appropriate
                result = z_test_proportions(int(x1), n1, int(x2), n2, alpha)
                recommendations.append("Use Z-test (normal approximation valid)")
                tests_run['z_test'] = result
            elif min_count >= 5:
                # Z-test with continuity correction or Fisher's
                result_z = z_test_proportions(int(x1), n1, int(x2), n2, alpha)
                table = np.array([[x1, n1-x1], [x2, n2-x2]])
                result_fisher = fishers_test_2x2(table, alpha)

                recommendations.append("Consider both Z-test (with caution) and Fisher's exact")
                tests_run['z_test'] = result_z
                tests_run['fisher'] = result_fisher
            else:
                # Fisher's exact test
                table = np.array([[x1, n1-x1], [x2, n2-x2]])
                result = fishers_test_2x2(table, alpha)
                recommendations.append("Use Fisher's exact test (small expected counts)")
                tests_run['fisher'] = result

    # Mean test
    elif metric_type == 'mean':
        if check_assumptions:
            # Check normality
            norm1 = check_normality(data1)
            norm2 = check_normality(data2)

            # Check variance equality
            var_eq = check_variance_equality(data1, data2, alpha)

            # Detect outliers
            outliers1 = detect_outliers(data1, method='iqr')
            outliers2 = detect_outliers(data2, method='iqr')

            normal1 = norm1['shapiro_normal']
            normal2 = norm2['shapiro_normal']
            equal_var = var_eq['levene_equal']
            has_outliers = (outliers1['n_outliers'] > 0 or
                           outliers2['n_outliers'] > 0)

            tests_run['normality_check'] = {
                'group1_normal': normal1,
                'group2_normal': normal2,
                'variance_equal': equal_var,
                'has_outliers': has_outliers
            }

            if normal1 and normal2:
                if equal_var:
                    result = t_test_means(data1, data2, alpha, equal_var=True)
                    recommendations.append("Use pooled t-test (normal data, equal variances)")
                else:
                    result = t_test_means(data1, data2, alpha, equal_var=False)
                    recommendations.append("Use Welch's t-test (normal data, unequal variances)")
                tests_run['t_test'] = result
            else:
                # Non-normal data
                if n1 >= 30 and n2 >= 30:
                    result_t = t_test_means(data1, data2, alpha, equal_var=False)
                    result_mw = mann_whitney_test(data1, data2, alpha)

                    recommendations.append("Large samples: t-test likely robust (CLT), "
                                         "but consider Mann-Whitney if heavily skewed")
                    tests_run['t_test'] = result_t
                    tests_run['mann_whitney'] = result_mw
                else:
                    result_mw = mann_whitney_test(data1, data2, alpha)
                    result_perm = permutation_test(data1, data2,
                                                   n_permutations=10000,
                                                   seed=42)

                    recommendations.append("Use Mann-Whitney U test or permutation test "
                                         "(non-normal, small samples)")
                    tests_run['mann_whitney'] = result_mw
                    tests_run['permutation'] = result_perm
        else:
            # Default: Welch's t-test (generally robust)
            result = t_test_means(data1, data2, alpha, equal_var=False)
            recommendations.append("Use Welch's t-test (default, robust choice)")
            tests_run['t_test'] = result

    # Ordinal or rank data
    elif metric_type == 'ordinal':
        result = mann_whitney_test(data1, data2, alpha)
        recommendations.append("Use Mann-Whitney U test (ordinal data)")
        tests_run['mann_whitney'] = result

    # Count data
    elif metric_type == 'count':
        # Poisson-based tests or transformations
        recommendations.append("Consider: (1) Square-root transform + t-test, "
                              "(2) Log transform + t-test, (3) Mann-Whitney, "
                              "(4) Generalized linear model")

        # Try multiple approaches
        result_t = t_test_means(data1, data2, alpha, equal_var=False)
        result_mw = mann_whitney_test(data1, data2, alpha)

        tests_run['t_test'] = result_t
        tests_run['mann_whitney'] = result_mw

    return {
        'metric_type': metric_type,
        'recommendations': recommendations,
        'tests_performed': list(tests_run.keys()),
        'results': tests_run,
        'sample_sizes': (n1, n2)
    }

# Example usage
print("=== Test Selection Framework ===\n")

np.random.seed(42)

# Example 1: Binary outcome (conversion rate)
conversions_control = np.random.binomial(1, 0.10, 2000)
conversions_treatment = np.random.binomial(1, 0.12, 2000)

print("Example 1: Conversion Rate")
result1 = select_statistical_test(conversions_control, conversions_treatment,
                                  metric_type='proportion')
print(f"Metric type: {result1['metric_type']}")
print(f"Recommendations:")
for rec in result1['recommendations']:
    print(f"  - {rec}")
print(f"Tests performed: {', '.join(result1['tests_performed'])}\n")

# Example 2: Continuous metric (normal data)
revenue_control = np.random.normal(50, 10, 500)
revenue_treatment = np.random.normal(53, 10, 500)

print("Example 2: Revenue (Normal Distribution)")
result2 = select_statistical_test(revenue_control, revenue_treatment,
                                  metric_type='mean', check_assumptions=True)
print(f"Recommendations:")
for rec in result2['recommendations']:
    print(f"  - {rec}")
if 'normality_check' in result2['results']:
    checks = result2['results']['normality_check']
    print(f"Assumption checks:")
    print(f"  Group 1 normal: {checks['group1_normal']}")
    print(f"  Group 2 normal: {checks['group2_normal']}")
    print(f"  Variances equal: {checks['variance_equal']}")
print()

# Example 3: Skewed data
pageviews_control = np.random.gamma(2, 5, 500)
pageviews_treatment = np.random.gamma(2.5, 5, 500)

print("Example 3: Page Views (Skewed Distribution)")
result3 = select_statistical_test(pageviews_control, pageviews_treatment,
                                  metric_type='mean', check_assumptions=True)
print(f"Recommendations:")
for rec in result3['recommendations']:
    print(f"  - {rec}")
\end{lstlisting}

\section{Common Pitfalls and How to Avoid Them}

\subsection{Pitfall 1: Peeking at Results}

\textbf{Problem:} Continuously checking p-values and stopping when significant inflates Type I error.

\textbf{Solution:} Pre-specify sample size and duration. If monitoring needed, use group sequential testing methods (Chapter 8).

\subsection{Pitfall 2: Multiple Testing Without Correction}

\textbf{Problem:} Testing many metrics without adjustment increases false positives.

\textbf{Solution:} Apply Bonferroni or FDR correction (Chapter 6). Pre-specify primary metric.

\subsection{Pitfall 3: Misinterpreting P-Values}

\textbf{Problem:} Confusing $P(\text{data}|H_0)$ with $P(H_0|\text{data})$.

\textbf{Solution:} Report effect sizes and confidence intervals. Focus on practical significance.

\subsection{Pitfall 4: Ignoring Assumptions}

\textbf{Problem:} Using t-tests on heavily skewed data with small samples.

\textbf{Solution:} Check assumptions visually and statistically. Use non-parametric alternatives when needed.

\subsection{Pitfall 5: Conflating Statistical and Practical Significance}

\textbf{Problem:} Declaring tiny effects "significant" with huge samples, or dismissing important effects as "non-significant" with small samples.

\textbf{Solution:} Always report effect sizes. Use minimum detectable effect (MDE) framework from Chapter 4.

\section{Summary}

Hypothesis testing provides a rigorous framework for A/B test analysis:

\begin{itemize}
    \item \textbf{Proportions:} Z-test for large samples, Fisher's exact for small
    \item \textbf{Means:} Welch's t-test as robust default
    \item \textbf{Categorical:} Chi-squared test (expected counts $\geq$ 5)
    \item \textbf{Assumptions matter:} Check normality, variance equality, independence
    \item \textbf{Non-parametric alternatives:} Mann-Whitney, permutation, bootstrap when assumptions fail
    \item \textbf{P-values are not enough:} Report effect sizes, confidence intervals, practical significance
\end{itemize}

Proper test selection and interpretation enable sound statistical inference and business decisions.

\section{Further Reading}

\begin{itemize}
    \item \cite{casella2002statistical}: Rigorous theoretical foundations
    \item \cite{good2013permutation}: Comprehensive treatment of permutation and bootstrap methods
    \item \cite{wasserstein2016asa}: ASA statement on p-values and statistical significance
\end{itemize}

\section{Exercises}

\subsection{Conceptual Questions}

\begin{enumerate}
    \item Explain why p-value is $P(\text{data}|H_0)$ and not $P(H_0|\text{data})$. What additional information would be needed to calculate $P(H_0|\text{data})$?

    \item When should you use Welch's t-test instead of pooled t-test? Why is Welch's often recommended as the default?

    \item Describe a scenario where statistical significance ($p < 0.05$) exists but practical significance does not, and vice versa.

    \item Why does "peeking" at A/B test results inflate Type I error? Sketch how error rate changes with number of peeks.
\end{enumerate}

\subsection{Calculation Exercises}

\begin{enumerate}
    \item Given: Control (n=1000, conversions=80), Treatment (n=1000, conversions=95)
    \begin{enumerate}
        \item Calculate z-statistic and p-value for two-sided test
        \item Construct 95% confidence interval for difference
        \item Calculate relative lift
        \item Verify $\chi^2 = Z^2$ using contingency table
    \end{enumerate}

    \item Given: Control (n=500, mean=\$45.20, sd=\$28.30), Treatment (n=500, mean=\$48.70, sd=\$31.10)
    \begin{enumerate}
        \item Perform Welch's t-test
        \item Calculate Cohen's d
        \item Construct 95% CI for difference
        \item Calculate variance ratio and assess assumption
    \end{enumerate}
\end{enumerate}

\subsection{Programming Exercises}

\begin{enumerate}
    \item Implement comprehensive test comparison framework:
    \begin{itemize}
        \item Generate data from various distributions (normal, skewed, heavy-tailed)
        \item Run t-test, Mann-Whitney, permutation test
        \item Compare p-values and power across methods
        \item Visualize results
    \end{itemize}

    \item Create assumption diagnostic dashboard:
    \begin{itemize}
        \item Automated normality checks
        \item Variance equality tests
        \item Outlier detection
        \item Visual diagnostics
        \item Test recommendations
    \end{itemize}

    \item Bootstrap confidence interval comparison:
    \begin{itemize}
        \item Implement percentile and BCa methods
        \item Compare to parametric CIs
        \item Evaluate coverage rates via simulation
        \item Test with various distributions
    \end{itemize}
\end{enumerate}

\subsection{Applied Exercises}

\begin{enumerate}
    \item \textbf{Real data analysis:} Obtain an A/B test dataset (public or from your organization)
    \begin{itemize}
        \item Perform comprehensive assumption checks
        \item Run appropriate tests
        \item Compare parametric and non-parametric methods
        \item Report results with effect sizes and CIs
        \item Make business recommendation
    \end{itemize}

    \item \textbf{Simulation study:} Investigate robustness of t-test
    \begin{itemize}
        \item Generate data from increasingly non-normal distributions
        \item Compare Type I error rates for t-test vs. Mann-Whitney
        \item Vary sample sizes
        \item Identify boundary where t-test breaks down
    \end{itemize}

    \item \textbf{Decision framework:} Create flowchart for test selection
    \begin{itemize}
        \item Decision points based on metric type, sample size, assumptions
        \item Implement as interactive tool
        \item Test with various scenarios
        \item Document edge cases
    \end{itemize}
\end{enumerate}

\subsection{Advanced Projects}

\begin{enumerate}
    \item \textbf{Meta-analysis:} Analyze 100+ A/B tests from an organization
    \begin{itemize}
        \item What proportion use appropriate tests?
        \item How often are assumptions checked?
        \item Do parametric and non-parametric yield different conclusions?
        \item Recommend process improvements
    \end{itemize}

    \item \textbf{Power simulation:} Compare power across test types
    \begin{itemize}
        \item Fix true effect size
        \item Vary distribution shape
        \item Calculate empirical power for each test
        \item Create power curves
        \item Identify optimal test for each scenario
    \end{itemize}
\end{enumerate}
